% Formato APA 7 para el informe. Se inyecta en el preambulo del documento.

% Imagenes.
\usepackage{graphicx}

% Margenes de 1 pulgada.
\usepackage{geometry}
\geometry{letterpaper, margin=1in, headheight=15pt}

% Interlineado 1.5.
\usepackage{setspace}
\onehalfspacing

% Sangria de 0.5 in en la primera linea, sin espacio entre parrafos.
\setlength{\parindent}{0.5in}
\setlength{\parskip}{0pt}

% Texto alineado a la izquierda conservando la sangria.
\usepackage{ragged2e}
\setlength{\RaggedRightParindent}{\parindent}
\RaggedRight

% Sin viudas ni huerfanas.
\widowpenalty=10000
\clubpenalty=10000
\displaywidowpenalty=10000
\brokenpenalty=10000
\raggedbottom

% Numero de pagina arriba a la derecha, sin linea.
\usepackage{fancyhdr}
\pagestyle{fancy}
\fancyhf{}
\fancyhead[R]{\thepage}
\renewcommand{\headrulewidth}{0pt}
\fancypagestyle{plain}{%
  \fancyhf{}\fancyhead[R]{\thepage}\renewcommand{\headrulewidth}{0pt}}

% Tamanos de letra arbitrarios y columnas con formato propio.
\usepackage{anyfontsize}
\usepackage{array}
\usepackage{xurl}

% Tablas con cuadricula completa. El filtro config/pdf-only.lua las escribe como
% longtable, por eso aqui se cargan los paquetes que necesita: longtable para
% que se partan entre paginas, multirow para las celdas combinadas y colortbl
% para el sombreado del encabezado.
\usepackage{longtable}
\usepackage{multirow}
\usepackage{colortbl}
\usepackage{etoolbox}

% Sin este parche una longtable que sigue a un encabezado de cuarto nivel se
% imprime antes que el propio encabezado.
\makeatletter
\patchcmd\longtable{\par}{\if@noskipsec\mbox{}\fi\par}{}{}
\makeatother

% Ancho util de cada tabla: lo que queda de \linewidth despues de descontar
% los separadores y las lineas verticales. El filtro reparte fracciones de esta
% longitud entre las columnas.
\newlength{\pdftblwidth}

% Las lineas de la cuadricula, algo mas finas que el trazo por defecto.
\setlength{\arrayrulewidth}{0.4pt}
\renewcommand{\arraystretch}{1.25}

% Los cinco niveles de encabezado de APA 7.
% Los numeros (1.1, 2.3.4) se escriben a mano en los .md, por eso no hay etiqueta.
\usepackage{titlesec}
\titleformat{\section}{\normalfont\bfseries\centering}{}{0pt}{}
\titleformat{\subsection}{\normalfont\bfseries\RaggedRight}{}{0pt}{}
\titleformat{\subsubsection}{\normalfont\bfseries\itshape\RaggedRight}{}{0pt}{}
\titleformat{\paragraph}[block]{\normalfont\bfseries}{}{0pt}{}
\titleformat{\subparagraph}[block]{\normalfont\bfseries\itshape}{}{0pt}{}
\titlespacing*{\section}{0pt}{2ex plus .5ex}{1ex}
\titlespacing*{\subsection}{0pt}{1.6ex plus .4ex}{0.8ex}
\titlespacing*{\subsubsection}{0pt}{1.4ex plus .4ex}{0.7ex}
\titlespacing*{\paragraph}{0pt}{1.2ex plus .3ex}{0.6ex}
\titlespacing*{\subparagraph}{0pt}{1.2ex plus .3ex}{0.6ex}

% Cada seccion de primer nivel (#) empieza en pagina nueva,
% como exige el enunciado para Registro de Versiones, Contenido, capitulos, etc.
\let\oldsection\section
\renewcommand{\section}{\clearpage\oldsection}

% El titulo del indice. Sin esto LaTeX lo rotula "Indice".
\AtBeginDocument{\renewcommand{\contentsname}{Contenido}}

% Figuras y tablas donde se escriben, con el rotulo arriba.
\usepackage{float}
\makeatletter\def\fps@figure{H}\def\fps@table{H}\makeatother
\usepackage{caption}
\captionsetup{position=above, labelfont=bf, textfont=it,
              singlelinecheck=false, justification=raggedright, skip=6pt}

% En espanol babel rotula "Cuadro"; APA usa "Tabla".
\AtBeginDocument{%
  \renewcommand{\tablename}{Tabla}%
  \renewcommand{\figurename}{Figura}}

% Sangria francesa de 0.5 in en las referencias.
% El \ifdefined evita fallar cuando se compila sin citas.
\AtBeginDocument{%
  \ifdefined\cslhangindent
    \setlength{\cslhangindent}{0.5in}%
    \setlength{\csllabelwidth}{0pt}%
  \fi}

% Pandoc pide cada imagen con height=\textheight, de modo que una figura alta se
% quedaba con la pagina entera y dejaba sola en la anterior la linea que la
% presenta. Al limitar la clave height de graphicx a 0.75 del alto del texto, el
% rotulo y la figura caben juntos y las imagenes pequenas no cambian.
\newlength{\pdfimageheight}
\makeatletter
\let\pdforiginalheightkey\KV@Gin@height
\define@key{Gin}{height}{%
  \setlength{\pdfimageheight}{#1}%
  \ifdim\pdfimageheight>0.75\textheight\setlength{\pdfimageheight}{0.75\textheight}\fi
  \pdforiginalheightkey{\the\pdfimageheight}}
\makeatother

% Bloques de codigo que se parten en vez de salirse del margen.
\usepackage{fvextra}
\DefineVerbatimEnvironment{Highlighting}{Verbatim}{%
  breaklines,breaknonspaceingroup,breakanywhere,commandchars=\\\{\}}
